\section{A concrete conjecture disproof, revalidated}\label{sec:density}

Let \(q\) be a prime power. A set \(E\subseteq\F_q[x]\setminus\{0\}\)
is polynomial-ratio progression-free if it contains no three elements
\[
 f,\quad rf,\quad r^2f
 \qquad(f\ne0,\quad r\in\F_q[x],\quad \deg r>0).
\]
The upper density is
\[
 \overline d(E)=\limsup_{N\to\infty}
       \frac{|\{f\in E:\deg f\le N\}|}{q^{N+1}}.
\]
Write \(\beta_q\) for the supremum of these upper densities and put
\[
 A=\{a\ge0:\text{every ternary digit of }a\text{ is }0\text{ or }1\},
 \qquad
 m_q=(1-q^{-1})\sum_{a\in A}q^{-a}.
\]
Since \(\sum_{a\in A}t^a=\prod_{j\ge0}(1+t^{3^j})\),
\begin{equation}\label{eq:mq}
 m_q=(1-q^{-2})\prod_{j\ge1}(1+q^{-3^j}).
\end{equation}
The original Conjecture~6.3 asserts \(\beta_q=m_q\)
\cite{asada}. It concerns polynomial ratios. Rational-function ratios are
a separate problem in that paper; the argument below does not solve it.

\begin{theorem}[Inherited ProveIt counterexample]\label{th:density}
For every prime power \(q\), there is one infinite polynomial-ratio
progression-free set \(E\subseteq\F_q[x]\) with
\[
 \overline d(E)>m_q.
\]
Consequently \(\beta_q>m_q\), contradicting the proposed value in
Conjecture~6.3.
\end{theorem}

The result and the construction are already in
\path{grh-methods-polynomial-density.zip} \cite{priorDensity}.
We supply the proof because a claimed resolution in a prior report should
be checked rather than repeated on its authority.

\subsection{The finite-replacement principle}

The ternary set \(A\) has no nonconstant three-term arithmetic progression.
Indeed, in an equality \(a+c=2b\) with all three in \(A\), neither side
has base-three carries. At each digit the only possibilities are
\(0+0=2\cdot0\) and \(1+1=2\cdot1\), so \(a=b=c\).

\begin{lemma}[Replacing degree deficits within rough-factor classes]
\label{lem:replacement}
Suppose \(D\subseteq\N\) is three-term arithmetic-progression-free and
\(D\setminus A\subseteq[0,L]\) for some integer \(L\ge1\).
Let \(\mathcal P\) be all monic irreducibles of degree at most \(L\), and let
\[
 H(u)=\prod_{P\in\mathcal P}(1-u^{\deg P})^{-1}
     =\sum_{n\ge0}H_nu^n.
\]
If, for some \(0<z<1\),
\[
 J(z):=\sum_{d\in D}z^d-\sum_{a\in A}z^a>0,
\]
then some integer \(R\ge0\) satisfies
\[
 \Delta_R=[u^R]H(u)J(u)\ge1.
\]
Moreover,
\begin{equation}\label{eq:delta}
 \beta_q\ge m_q+(q-1)\Delta_Rq^{-R-1}c_q,
 \qquad
 c_q=(1-q^{-1})\prod_{P\in\mathcal P}(1+q^{-\deg P})^{-1}>0.
\end{equation}
\end{lemma}

\begin{proof}
All series converge absolutely at \(z\). Their product has integer
coefficients, and \(H(z)J(z)>0\). If every coefficient were nonpositive,
the value at \(z\) would be nonpositive. Thus a coefficient
\(\Delta_R\ge1\) exists.

Call a monic polynomial \emph{smooth} if all its irreducible factors belong
to \(\mathcal P\), and \emph{rough} if none do. Every nonzero polynomial has
a unique expression
\[
 f=cgh,\qquad c\in\F_q^\times,\quad g\text{ monic rough},\quad
 h\text{ monic smooth}.
\]
At degree cutoff \(N\ge R\), the baseline construction includes the
polynomials with deficit \(d=N-\deg f\) in \(A\). It is progression-free
because degrees, and hence reversed deficits, would otherwise form a
nonconstant arithmetic progression.

For each squarefree rough monic \(g\) of degree \(N-R\), change the rule
on the entire class \(\{cgh\}\) to the deficit rule \(d\in D\).
Denote the resulting finite set by \(E_N\). A polynomial in this class
with deficit \(d\) has \(\deg h=R-d\). Consequently each \(g\) changes
the number of selected polynomials by exactly
\((q-1)\Delta_R\).

We check the delicate point: \emph{a new progression cannot escape the
modified class}. Any new progression contains an added polynomial
\(f=cgh\), whose deficit is in \(D\setminus A\), hence at most \(L\).
If this added polynomial is the first or middle term, a later term is
still inside the degree-\(N\) ball. Therefore the ratio has degree at most
\(L\), and every irreducible factor of the ratio is smooth. If the added
polynomial is the last term, the square of the ratio divides it.
Its rough part \(g\) is squarefree, so again every irreducible factor
of the ratio is smooth. In either case all three terms have the same
monic rough part \(g\), and their distinct deficits lie in \(D\).
That contradicts the progression-free property of \(D\).

Let \(G_n\) count squarefree rough monic polynomials of degree \(n\).
Unique factorization gives
\[
 \sum_{n\ge0}G_nu^n
   =\frac{1-qu^2}{1-qu}
      \prod_{P\in\mathcal P}(1+u^{\deg P})^{-1}.
\]
Subtracting its pole part at \(u=1/q\) leaves a rational function analytic
on \(|u|<1\). It follows by coefficient comparison that
\(G_n/q^n\to c_q\). The exact density of \(E_N\) is
\[
 \frac{|E_N|}{q^{N+1}}
  =(1-q^{-1})\sum_{\substack{a\in A\\a\le N}}q^{-a}
      +(q-1)\Delta_R\,\frac{G_{N-R}}{q^{N+1}}.
\]
It tends to the right side of \eqref{eq:delta}.

It remains to obtain a single infinite set, rather than a different set
for each cutoff. Choose rapidly increasing integers \(N_j\) such that
\[
 N_j-j>2N_{j-1},\qquad N_j-2j>N_{j-1}.
\]
From \(E_{N_j}\) keep only degrees in \([N_j-j,N_j]\), and take the
union over \(j\). A progression using two earlier terms has last degree
at most \(2N_{j-1}\). A progression using two terms in the latest block
has first degree at least \(N_j-2j\). Both possibilities contradict
the displayed separations. Progressions within one block were already
excluded. The deleted lower degrees have density at most \(q^{-j-1}\)
at cutoff \(N_j\), which tends to zero. This proves \eqref{eq:delta}
for one infinite set.
\end{proof}

\subsection{An explicit 74-point witness}

For the smallest inherited certificate used here, take \(D=B\), where
\[
\begin{split}
B=\{&
0,1,4,5,11,13,14,16,20,30,37,38,41,42,48,50,51,53,57,67,\\
&111,112,115,116,122,124,125,127,131,141,148,149,152,153,\\
&159,161,162,164,168,178,333,334,337,338,344,346,347,349,\\
&353,363,370,371,374,375,381,383,384,386,390,437,440,441,\\
&444,449,450,453,455,464,467,468,471,476,477,480\}.
\end{split}
\]
This is the translated \(878,\ldots,1358\) slice of Gavin Theobald's
\(a(143)\le1360\) record in Wroblewski's database, already selected by the
prior report \cite{wroblewski,priorDensity}. The present audit matched the
record against the primary database. More importantly, correctness
requires only the displayed list: all \(\binom{74}{2}=2701\) midpoint
tests were checked exactly. No assertion that this list is optimal is
used.

At \(z=127/128\), let
\[
 A_7=\left\{\sum_{j=0}^{6}\epsilon_j3^j:\epsilon_j\in\{0,1\}\right\}.
\]
It has 128 elements and maximum 1093. Clearing the denominator \(128^{1093}\)
proves the exact integer inequality
\begin{equation}\label{eq:finitegain}
 \sum_{b\in B}z^b-\sum_{a\in A_7}z^a>\frac{16}{1000}.
\end{equation}
Every omitted exponent in \(A\) is at least \(3^7=2187\), so
\[
 \sum_{a\in A\setminus A_7}z^a
 \le \frac{z^{2187}}{1-z}
 =128z^{2187}<\frac1{100000}.
\]
The last inequality was also checked by integer multiplication. Hence
\begin{equation}\label{eq:Jpositive}
 J(127/128)>\frac{3}{200}.
\end{equation}
Lemma~\ref{lem:replacement}, with \(L=480\), proves
Theorem~\ref{th:density}. The accompanying program
\path{code/check_density_certificate.py} verifies these exact inputs.
The asymptotic and infinite-set parts are the mathematical argument
above, not consequences of running a finite search.

\subsection{A database-free Behrend witness}

There is also a fully elementary alternative certificate, which displays
the connection with the Lean proof's combinatorial ingredient. Put
\(M=3^{20}\). For \(u\in\{0,1,\ldots,5\}^9\), define
\[
 b(u)=\sum_{i=0}^{8}u_i11^i,\qquad
 B_0=\left\{b(u):\sum_{i=0}^{8}(2u_i-5)^2=105\right\}.
\]
Each squared coordinate is \(1,9,\) or \(25\), with multiplicity two.
If the three counts are \(i,j,k\), then \(i+j+k=9\) and \(j+3k=12\).
Thus
\[
 |B_0|=2^9\left(
 \frac{9!}{1!\,6!\,2!}+\frac{9!}{3!\,3!\,3!}
                    +\frac{9!}{5!\,0!\,4!}\right)
 =1{,}053{,}696>2^{20}.
\]
No sum or doubled digit exceeds 10. A midpoint relation in base 11
therefore holds coordinatewise. Three points on a Euclidean sphere cannot
have one as the midpoint of the other two unless they coincide.
Consequently \(B_0\) is progression-free. Moreover,
\[
 \max B_0\le(11^9-1)/2<M/2.
\]

Set \(D=B_0\cup(A\cap[M,\infty))\). This remains progression-free.
Two low terms cannot reach a high term because \(2\max B_0<M\).
For a possible progression \(a<b<c\) with one low and two high terms,
let \(S=3^j\) be the largest ternary place of \(c\). If \(b<S\),
then \(2b\le S-1<c\). If \(b\ge S\), then
\[
 2b-c\ge(S+1)/2\ge(M+1)/2>a.
\]
Both alternatives exclude \(a+c=2b\). The difference
\[
 J(u)=\sum_{b\in B_0}u^b-\sum_{\substack{a\in A\\a<M}}u^a
\]
is a polynomial with \(J(1)=1{,}053{,}696-2^{20}>0\).
It is positive for some real \(z<1\) sufficiently close to 1, so
Lemma~\ref{lem:replacement} applies with \(L=M\).

The Lean source's \path{exists_three_term_free} uses the same carry-free
sphere idea to choose sparse embeddings in a CM-field construction
\cite{behrend,repo}. Here it selects polynomial degree deficits.
This is a transfer of a tool within the proof. The sphere method itself
predates the RH development by decades.

\subsection{What has and has not been resolved}

The strict inequality \(\beta_q>m_q\) answers the precise conjectured value
in Conjecture~6.3 in the negative. It does not determine \(\beta_q\), give a
large decimal improvement, or produce a natural-density result. The
construction deliberately uses separated blocks and a limsup. It also
does not treat arbitrary rational-function ratios, because the confinement
argument uses a polynomial ratio dividing the last term twice.

The proof has been reviewed in this session and its finite inputs were
rerun. It is not represented as independently published, peer reviewed,
or Lean formalized. Targeted searches did not locate a separate earlier
resolution, but the report does not assert a priority theorem.
